\subsection{可逆元素}

定义6. 设 E 是一个含幺广群， \(\top\) 为其合成律， \(e\) 为其单位元， \(x\) 和 \(x'\) 为 E 的两个元素。 \(x'\) 称为 \(x\) 的左逆（相应地，右逆，相应地，逆），如果 \(x' \top x = e\) （相应地 \(x \top x' = e\) ，相应地， \(x' \top x = x \top x' = e\) ).

E 的元素 x 称为左可逆（相应地，右可逆，相应地，可逆）的，如果它有左逆元（相应地，右逆元，相应地，逆元）。

一个所有元素都可逆的幺半群称为群。

“对称的”和“可对称化的”有时分别用来代替“逆”和“可逆”。当 E 上的律用加法记法书写时，我们通常说负元而不说逆元。

示例。(1) 单位元是其自身的逆元。

\begin{enumerate}
\def\labelenumi{(\arabic{enumi})}
\setcounter{enumi}{1}
\tightlist
\item
  在 E 到 E 的映射集合中，元素 \(f\) 是左可逆的（相应地，右可逆的），如果 \(f\) 是满射（相应地，单射）。左逆（相应地，右逆）就是与 \(f\) 相关联的收缩（相应地，截面）（集合论，II，§3，第8号，定义11）。 \(f\) 可逆的充要条件是 \(f\) 是一个双射。其唯一的逆就是 \(f\) 的逆双射.
\end{enumerate}

设 E 和 F 是两个含幺广群，f 是 E 到 F 的一个保单位同态。如果 \(x'\) 是 x 在 E 中的逆元，则 \(f(x')\) 是 \(f(x)\) 在 F 中的逆元。因此，若 x 是 E 的一个可逆元素，则 \(f(x)\) 是 F 的一个可逆元素。

特别地，若 R 是与含幺广群 E 上的律相容的等价关系，则 E 的可逆元素在 E/R 中的典范像也是可逆的。

命题3. 设 E 为一个幺半群，且 x 为 E 中的一个元素。

\begin{enumerate}
\def\labelenumi{(\roman{enumi})}
\item
  \(x\) 左（相应地，右）可逆的充要条件是由 \(x\) 所作的右（相应地，左）平移是满射。
\item
  \(x\) 可逆的充要条件是它既左可逆又右可逆。在这种情况下， \(x\) 有唯一的逆元，它同时也是其唯一的左（相应地，右）逆元。
\end{enumerate}

若 \(x'\) 是 \(x\) 的左逆，那么（命题1）

\[
\boldsymbol {\delta} _ {x} \circ \boldsymbol {\delta} _ {x ^ {\prime}} = \boldsymbol {\delta} _ {x ^ {\prime} \top x} = \boldsymbol {\delta} _ {e} = \mathrm{Id} _ {\mathrm{E}}
\]

且 \(\pmb{\delta}_{x}\) 是满射的。反之，如果 \(\pmb{\delta}_{x}\) 是满射的，则存在 E 的一个元素 \(x^{\prime}\) 使得 \(\pmb{\delta}_{x}(x^{\prime}) = e\) ，且 \(x^{\prime}\) 是 \(x\) 的左逆。(i) 的其他断言类似可得。

若 \(x'\) （相应地 \(x''\) ）是 \(x\) 的左（相应地，右）逆，则

\[
x ^ {\prime} = x ^ {\prime} \top e = x ^ {\prime} \top (x \top x ^ {\prime \prime}) = (x ^ {\prime} \top x) \top x ^ {\prime \prime} = e \top x ^ {\prime \prime} = x ^ {\prime \prime},
\]

由此得（ii）。

注。设 E 为一个幺半群，x 为 E 中的一个元素。若 x 左可逆，则 x 左可消；因为，若 \(x'\) 是 x 的左逆元，则

\[
\gamma_ {x ^ {\prime}} \circ \gamma_ {x} = \gamma_ {x ^ {\prime} \top x} = \gamma_ {e} = \mathrm{Id} _ {\mathrm{E}}
\]

且 \(\pmb{\gamma}_{x}\) 是单射的。特别地，如果 \(x\) 是左可逆的，则由 \(x\) 所作的左平移和右平移都是双射。反之，假设 \(\pmb{\gamma}_{x}\) 是双射；则存在 \(x' \in E\) 使得 \(xx' = \pmb{\gamma}_{x}(x') = e\) ；于是 \(\pmb{\gamma}_{x}(x' x) = (xx') x = x = \pmb{\gamma}_{x}(e)\) ，因此 \(x' x = e\) ，从而 \(x\) 是可逆的。类似地可以看出，如果 \(\delta_{x}\) 是双射，则 \(x\) 是可逆的。

命题4. 设 E 是一个幺半群，x 和 y 是 E 的可逆元素，其逆元分别为 \(x'\) 和 \(y'\) 。则 \(y' \top x'\) 是 \(x \top y\) 的逆.

这由以下关系得出

\[
\left(y ^ {\prime} \top x ^ {\prime}\right) \top (x \top y) = y ^ {\prime} \top \left(x ^ {\prime} \top x\right) \top y = y ^ {\prime} \top y = e
\]

以及对 \((x\top y)\top (y^{\prime}\top x^{\prime})\) 的类似计算

推论1. 设 E 为一个幺半群；若每个元素 \(x_{\alpha}\) （属于 E 的元素的一个有序序列 \((x_{\alpha})_{\alpha \in A}\) ）都有逆元 \(x_{\alpha}'\) ，则合成 \(\bigcap_{\alpha \in A} x_{\alpha}\) 有逆元 \(\bigcap_{\alpha \in A'} x_{\alpha}'\) ，其中 \(A'\) 是由 A 通过将其上的序替换为相反序而得到的全序集。

本推论可由命题4通过对A中元素个数进行归纳得出。

特别地，如果 \(x\) 和 \(x'\) 互为逆元，则 \(\frac{\pi}{\top} x\) 和 \(\frac{\pi}{\top} x'\) 对每个整数 \(n \geqslant 0\) 都互为逆元.

推论2. 在幺半群中，可逆元素的集合是稳定的。

命题5. 若在幺半群中 \(x\) 和 \(x'\) 互为逆元，且 \(x\) 与 \(y\) 可交换，则 \(x'\) 亦然。

由 \(x \top y = y \top x\) ，我们推出 \(x' \top (x \top y) \top x' = x' \top (y \top x) \top x'\) ，因此 \((x' \top x) \top (y \top x') = (x' \top y) \top (x \top x')\) ，即 \(y \top x' = x' \top y\) .

推论1. 设E为幺半群，X为E的子集，X'为X的中心化子。X'中每个可逆元素的逆元都属于X'。

推论2. 在幺半群中，中心可逆元素的逆元是中心元素。

